%% kapa: exam=1 round=34 year=2023 period=1 form=A subject=civil_law subject_ko=민법 seq=37 number=37
%% source: Q-Net 감정평가사 기출문제 (cst00309) artlSeq=5211789 file=raw/1008/5211789/2023년도 제34회 감정평가사 1차시험 1교시.pdf
%% conversion: kordoc 4.24.51 -> tools/gen1.py (auto); review=proofread-sonnet; flags=none
\begin{question}{37}
질권에 관한 설명으로 옳지 않은 것은? (다툼이 있으면 판례에 따름)
\begin{choices}
  \item 점유개정에 의한 동산질권설정은 인정되지 않는다.
  \item 질권자는 채권 전부를 변제받을 때까지 질물 전부에 대하여 그 권리를 행사할 수 있다.
  \item 질물이 공용징수된 경우, 질권자는 질권설정자가 받을 수용보상금에 대하여도 질권을 행사할 수 있다.
  \item 전질은 질물소유자인 질권설정자의 승낙이 있어도 허용되지 않는다.
  \item 부동산의 사용, 수익을 내용으로 하는 질권은 물권법정주의에 반한다.
\end{choices}
\end{question}
\answer{4}
